Dans le cadre de nos travaux, nous utilisons la librairie Jerboa. Cette
librairie, développée au laboratoire XLIM, permet la construction et la
manipulation de modèles géométriques à base topologique.

Les modèles 3D sont représentés à partir de leur structure topologique
(autrement dit leurs sommets, arêtes, faces, volumes et les relations entre ces
éléments) sous la forme de graphes topologiques. Ces derniers peuvent
être modifiés à l'aide de règles de transformations de graphes qui
permettent de filtrer une partie du graphe et de le remplacer par le
résultat de la transformation. Ces règles sont paramétrées, entre autres, par
des entités topologiques. Par exemple, la règle qui assemble deux volumes prend
en paramètre les deux faces le long desquelles les volumes seront cousus. Jerboa
permet de définir des opérations plus complexes construites par des scripts,
c'est-à-dire des enchaînements de règles.

Dans cette thèse, nous poursuivons l'approche de \cite{cardo2019}. Nous
considérons l'usage de brins comme entité topologique invariante. Les brins ne
peuvent être ni scindés ni fusionnés. Ils peuvent ainsi être utilisés pour
identifier les sommets, arêtes et faces d'un modèle.
%%% Local Variables:
%%% mode: latex
%%% TeX-master: "../main"
%%% End:
